\documentclass{article}
\usepackage[utf8]{inputenc}
\usepackage{geometry}
\geometry{a4paper, margin=1in}

\title{When Money and Power Don’t Teach}
\author{Albert Jan van Hoek}
\date{\today}

\begin{document}

\maketitle

If two people are in a relationship, and one consistently ends up with the money and the power, we usually call it toxic.

Yet when similar patterns appear in politics, we hesitate to name them.

\section*{The Vortex of Imbalance}

So imagine there are two people.

One has the money and the power.
The other does not.
The relationship is unequal.
How can such a relationship sustain itself?

\begin{itemize}
    \item Option 1. One person is starstruck and keeps giving money and power.
    \item Option 2. One person becomes narcissistic and keeps taking.
    \item Option 3. One person needs help. The other provides it — time, energy, money.
\end{itemize}

All three can slowly harden into the same thing: a self-sustaining vortex in which everyone is trying their best, yet the imbalance grows.

What makes this vortex feel inescapable is that money and power do not teach.

They accumulate, but they do not learn.

\section*{Relationships as Learning Systems}

But relationships — personal and political — are not just power structures.

They are interactions between learning systems.

And learning systems change when patterns become visible.
If both people can see the dynamic they are co-producing — not as blame, but as structure — something important happens: self-awareness enters the system.

Information can be shared. \\
Interpretations can be compared. \\
Behavior can adapt.

\section*{Theory of Long-Term Collaboration}

This is the core idea behind what I call the \textbf{Theory of Long-Term Collaboration}: when intelligent systems are able to reflect together on the relationship itself, collaboration becomes repairable.

Hope does not come from winning or redistributing power.
It comes from the fact that learning systems can update.

And love, in this sense, is not blind giving or endless endurance — it is the willingness to stay present, articulate the truth of the relationship, and trust that awareness can change behavior on both sides.

That possibility — fragile, demanding, but real — is what keeps relationships, and societies, alive.

\end{document}
